\documentclass[11pt]{article}

\usepackage[margin=25mm]{geometry}

\usepackage{amsmath,amssymb,fontspec,unicode-math}

\IfFontExistsTF{DejaVu Serif}{\setmainfont{DejaVu Serif}}{\setmainfont{Latin Modern Roman}}

\IfFontExistsTF{Latin Modern Math}{\setmathfont{Latin Modern Math}}{\setmathfont{DejaVu Math TeX Gyre}}

\usepackage{xurl,hyperref}

\setlength{\parindent}{0pt}\setlength{\parskip}{6pt}

\begin{document}

\title{MultiNet: Adaptive Multi-Viewed Subgraph Convolutional Networks for Graph Classification}

\author{ELVIA preliminary AI analysis}\date{21 October 2026}\maketitle

Preliminary assessment: Major.\par

Independent scientific review remains pending. Author adoption has not been requested.

The stationary-limit formula and spectral-radius argument do not establish MultiNet's central anti-collapse explanation. Explicit linear counterexamples show the failures. Other implementation, cross-validation and reporting discrepancies remain local findings; they do not establish which private runs produced the reported classifier scores.

\section*{Review scope}

Claim-focused reading of the identified paper and author-released materials. The arguments below distinguish mathematical contradictions, version-specific implementation defects and unresolved empirical consequences. No author code was executed in this revision.

\section*{Problem 1: Correct stationary limit and a graph-level counterexample}

Source: Main PDF p.4, Section 2.2, Theorem 1, Eqs 1-2 Main PDF p.5, Section 3.2

Assessment: supported; impact: major

\subsection*{Source claim (paraphrase)}

Theorem 1 identifies the deep symmetric-normalized propagation limit with an arithmetic-mean projector and uses it to infer graph-level collapse.

\subsection*{Analysis}

For a connected aperiodic graph with augmented degrees \(\widetilde d_i\), the normalized adjacency power converges to \(uu^\top\), where \(u_i=\sqrt{\widetilde d_i}/\sqrt{\sum_j\widetilde d_j}\). For the self-looped three-node path, \(u=(\sqrt2,\sqrt3,\sqrt2)/\sqrt7\), which is not the uniform vector. Thus \(uu^\top\ne\mathbf1\mathbf1^\top/3\). In the paper's linear identity-weight setting, equal initial features on this path and a regular triangle need not have equal limiting mean readouts: degree factors remain. The arithmetic-mean limit and the resulting universal graph-level conclusion therefore fail.

With constant scalar initial node feature \(1\), the self-looped path has limiting mean readout \((11+4\sqrt6)/21\), while the self-looped regular triangle has limiting mean \(1\). Their equal initial feature means therefore do not imply equal limiting graph readouts under the theorem’s own linear identity-weight simplification.

\subsection*{Finding and scope}

The counterexample refutes the stated stationary identity and graph-level implication. It does not deny empirical oversmoothing in graph networks.

\subsection*{Proposed correction}

Use the degree-dependent projector and specify conditions on limiting graph readouts. Distinguish mean/sum pooling, regular/disconnected graphs, weights and nonlinearities.

Proposed correction; independent verification and author adoption remain unestablished.

\section*{Problem 2: Smaller spectral radius does not establish anti-collapse}

Source: Main PDF p.7, Section 4, Eqs 11-12 and following paragraph

Assessment: supported; impact: major

\subsection*{Source claim (paraphrase)}

Section 4 treats a reduced spectral radius as sufficient evidence against collapse.

\subsection*{Analysis}

Admissible equal two-view gates produce \(M=S/2\), with spectral radius \(1/2\). If \(S^L\to uu^\top\), then \(M^L=2^{-L}S^L\to0\): features collapse to zero. Relative eigenvalue ratios are unchanged by this scalar rescaling, so normalized oversmoothing is not removed either.

\subsection*{Finding and scope}

Reducing the absolute spectral radius does not establish the claimed anti-collapse property. Learned gates and nonlinearities need their own analysis; classifier accuracy is not determined by this linear example.

\subsection*{Proposed correction}

Analyze relative spectral decay under the actual recurrence. A qualified linear remedy is persistent initial-feature injection with convergent resolvent; its assumptions and effect on accuracy would still require verification.

Proposed correction; independent verification and author adoption remain unestablished.

\section*{Problem 3: Evaluation-time BatchNorm and node-order dependence}

Source: code/model.py:118-120,125,134

Assessment: supported; impact: minor

\subsection*{Source claim (paraphrase)}

The graph classifier's evaluation should respect node relabeling and its stated normalization convention.

\subsection*{Analysis}

Fresh BatchNorm1d uses node index as channel and remains training-mode even after parent.eval(); it is unregistered and called unconditionally once. Helper batch-change error \(2.15487\); five full-model seeds produce relabeling differences \(0.000326\)–\(0.007388\).

\subsection*{Finding and scope}

Symmetric full-model input yielded only \(1.19\times10^{-7}\) relabeling difference; that negative result is retained. Synthetic differences are not benchmark accuracy changes.

\subsection*{Proposed correction}

Use registered, appropriately masked feature normalization or per-node LayerNorm; preserve evaluation mode and node equivariance.

Proposed correction; independent verification and author adoption remain unestablished.

\section*{Problem 4: Parameter registration can be lost on device transfer}

Source: code/model.py:20-22,71-73

Assessment: supported; impact: minor

\subsection*{Source claim (paraphrase)}

The learned convolution parameters must remain registered after device placement.

\subsection*{Analysis}

nn.Parameter(cpu\_tensor).to(device) returns an ordinary non-leaf Tensor when a device copy occurs, losing module parameter registration. Actual GraphConv has weight/bias on CPU but no named parameters on meta device.

\subsection*{Finding and scope}

CUDA execution was not performed. A no-op CPU conversion can conceal the issue; no inference about undisclosed historical training is made.

\subsection*{Proposed correction}

Create tensors on the target device before wrapping, or move the completed module. Check named\_parameters, state\_dict and optimizer membership; retrain after repair.

Proposed correction; independent verification and author adoption remain unestablished.

\section*{Problem 5: Padding is treated as graph content}

Source: code/main.py:104-126 code/model.py:134-135,173-177

Assessment: supported; impact: minor

\subsection*{Source claim (paraphrase)}

Zero padding should not become extra graph content in normalization or readout.

\subsection*{Analysis}

Real-node counts are not used as a mask; affine/normalization/readout operations include padded positions. Adding two zero-padded nodes to an unchanged graph changes full-model logits by 0.137769.

\subsection*{Finding and scope}

Synthetic counterexample establishes release sensitivity, not its effect under the authors fixed dataset padding. Corrected model accuracy is untested.

\subsection*{Proposed correction}

Pass real-node masks; zero padded features after transformations and exclude padding from assignments, pooling and normalization.

Proposed correction; independent verification and author adoption remain unestablished.

\section*{Problem 6: The implemented recurrence differs from the analyzed recurrence}

Source: Main PDF pp.5-6, Eqs 4-7 code/model.py:122-136,166-172

Assessment: supported; impact: minor

\subsection*{Source claim (paraphrase)}

The released recurrence is intended to implement the multiview propagation analyzed in the manuscript.

\subsection*{Analysis}

Paper specifies fixed initial assignments and separate view states propagated before final fusion. Release recomputes assignments and fuses every layer. Literal Eqs \(5\)-\(6\) also gate the first state twice, while release gates once; release adds normalization operations.

\subsection*{Finding and scope}

Alternative implemented architecture can be legitimate; discrepancy does not identify which private code generated tables.

\subsection*{Proposed correction}

Choose a canonical recurrence: preserve view states/fixed assignments as written, or revise equations and analyze layer-indexed assignments/fusion. Publish actual run configuration.

Proposed correction; independent verification and author adoption remain unestablished.

\section*{Problem 7: Cross-validation and coverage are not reproduced by the entrypoint}

Source: Main PDF p.8, Section 5.1 code/main.py:146-166,223-240

Assessment: supported; impact: minor

\subsection*{Source claim (paraphrase)}

The manuscript describes repeated ten-fold cross-validation for its graph-classification results.

\subsection*{Analysis}

Entrypoint uses fold \(0\) once rather than looping over ten repetitions of ten-fold CV. With \(188\) graphs, floor-width folds test only \(180\) indices; \(180\)-\(187\) are never tested and excluded from final-fold training.

\subsection*{Finding and scope}

Public reproduction path is incomplete; this does not prove the authors omitted their stated protocol in private scripts.

\subsection*{Proposed correction}

Use stratified complete outer folds and inner selection; preserve splits/seeds and publish \(100\) outer-fold scores plus repetition summaries.

Proposed correction; independent verification and author adoption remain unestablished.

\section*{Problem 8: Additional concrete reproduction faults}

Source: code/main.py:76,108-109,114-115,195-200,235-238 CLI and evaluation paths in code/main.py

Assessment: supported; impact: minor

\subsection*{Source claim (paraphrase)}

The released command-line and evaluation interfaces are intended to support the listed dataset settings.

\subsection*{Analysis}

Two-class head cannot accept six-class targets; two dimension arguments become strings when explicitly supplied. Final returned model is epoch \(99\) but accuracy is epoch \(90\); graph truncation uses first node indices and squeeze can scalarize singleton batches.

\subsection*{Finding and scope}

No benchmark rerun or corrected accuracy. Truncation is a defined release behavior requiring an explicit dataset protocol.

\subsection*{Proposed correction}

Derive class count, type integer arguments, return matching checkpoint/score, support dataset capacity or sparse batching, and use reshape(\(-1\)) for labels.

Proposed correction; independent verification and author adoption remain unestablished.

\section*{Problem 9: Error bars have incompatible definitions}

Source: Main PDF p.8, Section 5.1 and Table 1 Main PDF p.10, Table 3 Main PDF p.16, checklist item 7

Assessment: supported; impact: minor

\subsection*{Source claim (paraphrase)}

The uncertainty annotations describe the same statistic and independent evaluation unit.

\subsection*{Analysis}

Text describes standard deviation while captions/comparison/checklist describe standard error. Imported baseline rows may use incompatible uncertainty conventions; fold scores are not independent repetition-level measurements.

\subsection*{Finding and scope}

No independent significance claim is verified; the actual historical uncertainty calculation is unknown.

\subsection*{Proposed correction}

Publish fold/repetition scores, identify independent unit and each plus/minus quantity, and recompute consistent paired summaries.

Proposed correction; independent verification and author adoption remain unestablished.

\section*{Problem 10: Average-rank audit}

Source: Main PDF p.8, Table 1

Assessment: supported; impact: minor

\subsection*{Source claim (paraphrase)}

Table 1's average ranks summarize the stated within-dataset ranking convention.

\subsection*{Analysis}

Standard descending within-dataset ranks with average ties give DGK \(12.5714\) vs \(10.57\), RWGK \(14\) vs \(13.14\), WLSK \(8.1429\) vs 7.14. MultiNet mean rank \(16/7\simeq2.285714\) is reproduced and remains strongest under this convention.

\subsection*{Finding and scope}

An unstated alternative ranking convention could explain differences; classified unresolved rather than a proved historical calculation error.

\subsection*{Proposed correction}

Publish ranking universe, code and handling of missing entries; avoid mixing imported ranks from different tables.

Proposed correction; independent verification and author adoption remain unestablished.

\section*{Problem 11: The directional-diversity proxy is explicitly distinguished from the formal definition}

Source: Main PDF pp.3-4, Definition 1 Main PDF p.9, Eq 13 and Section 5.3

Assessment: not\_issue; impact: no\_issue\_found

\subsection*{Source claim (paraphrase)}

Directional diversity is compared with the formal diversity definition.

\subsection*{Analysis}

Authors acknowledge cosine-distance measure does not satisfy Definition 1. Vectors \((1,0)\) and \((2,0)\) have zero cosine distance but unequal magnitude; zero vectors need a convention.

\subsection*{Finding and scope}

The authors acknowledge this proxy limitation. It is a scoped clarification, not a current author error or a newly established anti-collapse guarantee.

\section*{Problem 12: Claim calibration and checklist consistency}

Source: Main conclusion Main PDF p.8, Table 1 Limitations and existing-assets checklist items Section 5.4

Assessment: supported; impact: minor

\subsection*{Source claim (paraphrase)}

The conclusion and checklist should reflect the displayed dataset wins, assumptions and assets.

\subsection*{Analysis}

Table 1 has MultiNet losses on ENZYMES and DD, so universal superiority exceeds five wins out of seven. Checklist NA answers conflict with disclosed assumptions/space limitations and use of public datasets; three variants are said but four listed.

\subsection*{Finding and scope}

Reporting corrections do not allege license violations or independently invalidate experiments.

\subsection*{Proposed correction}

State five dataset wins and strong average rank; disclose limitations/assets and correct variant count.

Proposed correction; independent verification and author adoption remain unestablished.

\section*{Limits}

Independent scientific adjudication remains pending.\par

No author code, model training or full benchmark evaluation was executed in this revision.\par

Version-specific source defects do not establish which code produced the published experiments.\par

Corrected benchmark scores, empirical effect sizes and author adoption remain unverified unless explicitly stated.\par

\section*{Source identity}

Current official NeurIPS 2025 proceedings PDF, freshly retrieved 21 October 2026

\url{https://proceedings.neurips.cc/paper_files/paper/2025/file/619cbddb92b8c6fecaf2b86463153be9-Paper-Conference.pdf}

Reviewed PDF SHA-256: \nolinkurl{1aeb4903475a5094c6da9ecb8b84bdade301528b693a37466e9d870d992cfa7e}

This current primary is the same research work as the predecessor, identified by matching title, authors and technical content, on the stated current publication channel. Its exact inspected bytes are identified by SHA-256; the earlier immutable source identity remains in the predecessor record. The download timestamp is recorded from this fresh primary-source retrieval.

\end{document}
